\chapter{2012年福建卷（理）}

\section{单选题}

\begin{problem}
若复数 \(z\) 满足 \(z\i = 1 - \i\)，则 \(z\) 等于（\quad）

\choices
    {\(-1 - \i\)}
    {\(1 - \i\)}
    {\(-1 + \i\)}
    {\(1 + \i\)}
\end{problem}

\begin{answer}
A
\end{answer}

\begin{solution}
设 \(z=a+b\i\)，其中 \(a\)，\(b\in\R\). 则
\[
z\i=(a+b\i)\i=-b+a\i=1-\i
\]
比较实部与虚部，得 \(-b=1\)，\(a=-1\)，所以
\[
z=-1-\i
\]
故选择 A.
\end{solution}

\begin{problem}
等差数列 \(\{a_{n}\}\) 中，\(a_{1} + a_{5} = 10\)，\(a_{4} = 7\)，则数列 \(\{a_{n}\}\) 的公差为（\quad）

\choices
    {\(1\)}
    {\(2\)}
    {\(3\)}
    {\(4\)}
\end{problem}

\begin{answer}
B
\end{answer}

\begin{solution}
设等差数列的公差为 \(d\). 由等差数列的性质，
\[
a_{1}+a_{5}=2a_{3}=10
\]
所以 \(a_{3}=5\). 又 \(a_{4}=a_{3}+d=7\)，故
\[
d=7-5=2
\]
因此选择 B.
\end{solution}

\begin{problem}
下列命题中，真命题是（\quad）

\begin{circlelist}
\item \(\exists x_{0} \in \R\)，\(\e^{x_{0}} \leqslant 0\)；
\item \(\forall x \in \R\)，\(2^{x} > x^{2}\)；
\item \(a + b = 0\) 的充要条件是 \(\frac{a}{b} = -1\)；
\item \(a > 1\)，\(b > 1\) 是 \(ab > 1\) 的充分条件.
\end{circlelist}

\choices
    {\(\circled{1}\circled{2}\)}
    {\(\circled{1}\circled{3}\)}
    {\(\circled{2}\circled{4}\)}
    {\(\circled{3}\circled{4}\)}
\end{problem}

\begin{answer}
D
\end{answer}

\begin{solution}
逐项判断. 因为对任意 \(x\in\R\) 都有 \(\e^{x}>0\)，命题 \(\circled{1}\) 错误. 取 \(x=2\)，有 \(2^{x}=x^{2}=4\)，不满足严格不等式，故命题 \(\circled{2}\) 错误. 命题 \(\circled{3}\) 需要先保证 \(b\neq0\)；当 \(a=b=0\) 时 \(a+b=0\) 成立，但 \(a/b\) 无意义，故命题 \(\circled{3}\) 错误. 若 \(a>1\)，\(b>1\)，则 \(ab>a>1\)，所以命题 \(\circled{4}\) 正确.

因此真命题只有 \(\circled{4}\)，选择 D.
\end{solution}

\begin{problem}
一个几何体的三视图形状都相同，大小均相等，那么这个几何体不可以是（\quad）

\choices
    {球}
    {三棱锥}
    {正方体}
    {圆柱}
\end{problem}

\begin{answer}
D
\end{answer}

\begin{solution}
球的三个投影都是圆，正方体的三个投影都是全等正方形. 对于正四面体，可将其四个顶点置于
\((1,1,1)\)、\((1,-1,-1)\)、\((-1,1,-1)\)、\((-1,-1,1)\). 此时分别向三个坐标平面作正投影，三个投影的顶点集合都为 \((\pm1,\pm1)\)，因而都是全等正方形. 圆柱的俯视图是圆，而正视图、侧视图是矩形，圆与矩形不可能形状相同且大小相等. 因此圆柱不符合条件，选择 D.
\end{solution}

\begin{problem}
下列不等式一定成立的是（\quad）

\choices
    {\(\lg\left(x^{2}+\frac{1}{4}\right) > \lg x\)（\(x > 0\)）}
    {\(\sin x + \frac{1}{\sin x} \geqslant 2\)（\(x \neq k\pi\)，\(k \in \Z\)）}
    {\(x^{2} + 1 \geqslant 2|x|\)（\(x \in \R\)）}
    {\(\frac{1}{x^{2} + 1} > 1\)（\(x \in \R\)）}
\end{problem}

\begin{answer}
C
\end{answer}

\begin{solution}
选项 A 中，由对数函数的单调性，需有
\(x^{2}+\frac14>x\)，\((x>0)\)
但 \(x^{2}-x+\frac14=(x-\frac12)^{2}\)，当 \(x=\frac12\) 时等号成立，所以 A 不一定成立.

选项 B 中，取 \(x=\frac{3\pi}{2}\)，则 \(\sin x=-1\)，从而
\[
\sin x+\frac{1}{\sin x}=-2<2
\]
故 B 错误.

选项 C 中，
\[
x^{2}+1-2|x|=(|x|-1)^{2}\geqslant0
\]
所以 C 对任意 \(x\in\R\) 都成立. 选项 D 在 \(x=0\) 时变为 \(1>1\)，错误.

因此选择 C.
\end{solution}

\begin{problem}
如图所示，在边长为 \(1\) 的正方形 \(OABC\) 中任取一点 \(P\)，则点 \(P\) 恰好取自阴影部分的概率为（\quad）

\bitmapfigure[max width=0.42\linewidth]{img/2012/fujian_science/q6_fig1.png}

\choices
    {\(\frac{1}{4}\)}
    {\(\frac{1}{5}\)}
    {\(\frac{1}{6}\)}
    {\(\frac{1}{7}\)}
\end{problem}

\begin{answer}
C
\end{answer}

\begin{solution}
正方形面积为 \(1\). 由图，阴影部分是曲线 \(y=\sqrt{x}\) 与直线 \(y=x\) 在 \(0\leqslant x\leqslant1\) 之间的区域，因此其面积为
\[
S=\int_{0}^{1}(\sqrt{x}-x)\,\mathrm{d}x
=\left.\frac23x^{3/2}-\frac12x^{2}\right|_{0}^{1}
=\frac23-\frac12=\frac16
\]
点在正方形内均匀取值，所求概率等于面积之比，故
\[
P=\frac{S}{1}=\frac16
\]
因此选择 C.
\end{solution}

\begin{problem}
设函数 \(D(x) = \begin{cases}
1, & x\text{ 为有理数},\\
0, & x\text{ 为无理数},
\end{cases}\) 则下列结论错误的是（\quad）

\choices
    {\(D(x)\) 的值域为 \(\{0,1\}\)}
    {\(D(x)\) 是偶函数}
    {\(D(x)\) 不是周期函数}
    {\(D(x)\) 不是单调函数}
\end{problem}

\begin{answer}
C
\end{answer}

\begin{solution}
有理数的相反数仍为有理数，无理数的相反数仍为无理数，所以 \(D(-x)=D(x)\)，命题 \(\circled{2}\) 正确. 取有理数和无理数分别作为自变量，函数值分别为 \(1\) 和 \(0\)，故值域为 \(\{0,1\}\)，命题 \(\circled{1}\) 正确.

对任意 \(x\in\R\)，\(x\) 与 \(x+1\) 同时为有理数或同时为无理数，所以 \(D(x+1)=D(x)\)，即 \(1\) 是其一个周期，命题 \(\circled{3}\) “不是周期函数”错误. 为说明它不是单调函数，取
\(\frac12<\sqrt2\)，有 \(D(\frac12)=1>D(\sqrt2)=0\)，故它不是单调递增函数；又取 \(\sqrt2<\frac32\)，有 \(D(\sqrt2)=0<D(\frac32)=1\)，故它不是单调递减函数. 因此命题 \(\circled{4}\) 正确.

因此错误的结论是 \(\circled{3}\)，选择 C.
\end{solution}

\begin{problem}
已知双曲线 \(\frac{x^{2}}{4} - \frac{y^{2}}{b^{2}} = 1\) 的右焦点与抛物线 \(y^{2} = 12x\) 的焦点重合，则该双曲线的焦点到其渐近线的距离等于（\quad）

\choices
    {\(\sqrt{5}\)}
    {\(4\sqrt{2}\)}
    {3}
    {5}
\end{problem}

\begin{answer}
A
\end{answer}

\begin{solution}
抛物线 \(y^{2}=12x\) 可写为 \(y^{2}=4px\)，故 \(p=3\)，其焦点为 \((3,0)\). 双曲线的实半轴为 \(a=2\)，右焦点为 \((c,0)\)，由焦点重合得 \(c=3\). 于是
\[
b^{2}=c^{2}-a^{2}=9-4=5
\]
双曲线的渐近线为
\[
y=\pm\frac{b}{a}x=\pm\frac{\sqrt5}{2}x
\]
取焦点 \(F(3,0)\)，它到任一渐近线的距离为
\[
d=\frac{\frac{3\sqrt5}{2}}{\sqrt{1+\frac54}}=\sqrt5
\]
故选择 A.
\end{solution}

\begin{problem}
若函数 \(y = 2^{x}\) 图象上存在点 \((x,y)\) 满足约束条件 \(\begin{cases}
x + y - 3 \leqslant 0,\\
x - 2y - 3 \leqslant 0,\\
x \geqslant m,
\end{cases}\) 则实数 \(m\) 的最大值为（\quad）

\choices
    {\(\frac{1}{2}\)}
    {1}
    {\(\frac{3}{2}\)}
    {2}
\end{problem}

\begin{answer}
B
\end{answer}

\begin{solution}
在函数图象上有 \(y=2^{x}\). 第一个约束变为
\[
g(x)=x+2^{x}-3\leqslant0
\]
因为 \(g'(x)=1+2^{x}\ln2>0\)，所以 \(g(x)\) 在 \(\R\) 上递增；又 \(g(1)=1+2-3=0\)，故第一个约束等价于 \(x\leqslant1\). 对于这样的 \(x\)，第二个约束满足
\[
x-2y-3=x-2^{x+1}-3<x-3\leqslant-2<0
\]
因而自动成立.

所以满足前两个约束的点对应的 \(x\) 的取值范围为 \((-\infty,1]\). 再结合 \(x\geqslant m\)，存在满足条件的点当且仅当 \(m\leqslant1\)，故 \(m\) 的最大值为 \(1\)，选择 B.
\end{solution}

\begin{problem}
函数 \(f(x)\) 在 \([a,b]\) 上有定义，若对任意 \(x_{1}\)，\(x_{2} \in [a,b]\)，有 \(f\left(\frac{x_{1} + x_{2}}{2}\right) \leqslant \frac{1}{2}[f(x_{1}) + f(x_{2})]\)，则称 \(f(x)\) 在 \([a,b]\) 上具有性质 \(P\). 设 \(f(x)\) 在 \([1,3]\) 上具有性质 \(P\)，现给出如下命题：

\begin{circlelist}
\item \(f(x)\) 在 \([1,3]\) 上的图象是连续不断的；
\item \(f(x^{2})\) 在 \([1,\sqrt{3}]\) 上具有性质 \(P\)；
\item 若 \(f(x)\) 在 \(x = 2\) 处取得最大值 \(1\)，则 \(f(x) = 1\)，\(x \in [1,3]\)；
\item 对任意 \(x_{1}\)，\(x_{2}\)，\(x_{3}\)，\(x_{4} \in [1,3]\)，有 \(f\left(\frac{x_{1} + x_{2} + x_{3} + x_{4}}{4}\right) \leqslant \frac{1}{4}[f(x_{1}) + f(x_{2}) + f(x_{3}) + f(x_{4})]\).
\end{circlelist}

其中真命题的序号是（\quad）

\choices
    {\(\circled{1}\circled{2}\)}
    {\(\circled{1}\circled{3}\)}
    {\(\circled{2}\circled{4}\)}
    {\(\circled{3}\circled{4}\)}
\end{problem}

\begin{answer}
D
\end{answer}

\begin{solution}
命题 \(\circled{1}\) 不成立. 例如在 \([1,3]\) 上定义
\[
f(x)=\begin{cases}
1,&x=1,\\
0,&1<x\leqslant3
\end{cases}
\]
当两个自变量都为 \(1\) 时性质 \(P\) 显然成立；当至少一个自变量大于 \(1\) 时，中点大于 \(1\)，左端为 \(0\)，右端非负，因此性质 \(P\) 仍成立. 但该函数在 \(x=1\) 处不连续，所以 \(\circled{1}\) 错误.

命题 \(\circled{2}\) 不成立. 取 \(f(x)=-x\)，它在 \([1,3]\) 上满足性质 \(P\)，因为两端平均值等于中点函数值. 此时 \(g(x)=f(x^{2})=-x^{2}\). 取 \(x_{1}=1\)，\(x_{2}=\sqrt3\)，则
\[
g\left(\frac{x_{1}+x_{2}}{2}\right)=-1-\frac{\sqrt3}{2}>-2=\frac{g(x_{1})+g(x_{2})}{2}
\]
不满足性质 \(P\)，所以 \(\circled{2}\) 错误.

对于命题 \(\circled{3}\)，任取 \(x\in[1,3]\)，令 \(y=4-x\)，则 \(y\in[1,3]\)，且 \((x+y)/2=2\). 由性质 \(P\) 及 \(f(2)=1\) 为最大值，
\[
1=f(2)\leqslant\frac{f(x)+f(4-x)}{2}\leqslant1
\]
两边只能处处取等号，故 \(f(x)=f(4-x)=1\)，命题 \(\circled{3}\) 正确.

对于命题 \(\circled{4}\)，先对 \(x_{1}\)，\(x_{2}\) 及 \(x_{3}\)，\(x_{4}\) 使用性质 \(P\)，再对两个中点使用一次性质 \(P\)，得
\[
f\left(\frac{x_{1}+x_{2}+x_{3}+x_{4}}{4}\right)
\leqslant\frac12\left[f\left(\frac{x_{1}+x_{2}}2\right)+f\left(\frac{x_{3}+x_{4}}2\right)\right]
\leqslant\frac14\left[f(x_{1})+f(x_{2})+f(x_{3})+f(x_{4})\right]
\]
所以 \(\circled{4}\) 正确. 真命题为 \(\circled{3}\)、\(\circled{4}\)，选择 D.
\end{solution}

\section{填空题}

\begin{problem}
\((a + x)^{4}\) 的展开式中 \(x^{3}\) 的系数等于 \(8\)，则实数 \(a =\)\fillinblank{}.
\end{problem}

\begin{answer}
\(2\)
\end{answer}

\begin{solution}
由二项式定理，\((a+x)^{4}\) 中含 \(x^{3}\) 的项为
\[
\mathrm C_4^3a x^{3}=4a x^{3}
\]
题设给出 \(4a=8\)，所以 \(a=2\).
\end{solution}

\begin{problem}
阅读如图所示的程序框图，运行相应的程序，输出的 \(s\) 值等于\fillinblank{}.

\begin{center}
\texfigure[float=inline,max width=0.44\linewidth]{img_repaint/2012/fujian_science/q12_fig1.tex}
\end{center}
\end{problem}

\FloatBarrier

\begin{answer}
\(-3\)
\end{answer}

\begin{solution}
流程图先判断 \(k<4\)，判断为“是”后依次执行 \(s=2s-k\) 和 \(k=k+1\). 初始值为 \(k=1\)，\(s=1\). 各次循环为
\begin{center}
\begin{tabular}{c|c|c|c}
\text{循环前 }\(k\)&\text{循环前 }\(s\)&\text{更新后 }\(s=2s-k\)&\text{更新后 }\(k=k+1\)\\ \hline
1&1&1&2\\
2&1&0&3\\
3&0&\(-3\)&4
\end{tabular}
\end{center}
此时 \(k=4\)，判断 \(k<4\) 为“否”，程序输出 \(s=-3\).
\end{solution}

\begin{problem}
已知 \(\triangle ABC\) 的三边长成公比为 \(\sqrt{2}\) 的等比数列，则其最大角的余弦值为\fillinblank{}.
\end{problem}

\begin{answer}
\(-\frac{\sqrt{2}}{4}\)
\end{answer}

\begin{solution}
设等比数列的三项依次为 \(t\)，\(t\sqrt2\)，\(2t\)，其中 \(t>0\). 最大边为 \(2t\)，设其所对的最大角为 \(C\). 由余弦定理，
\[
\cos C=\frac{t^{2}+(t\sqrt2)^{2}-(2t)^{2}}{2\cdot t\cdot t\sqrt2}
=\frac{t^{2}+2t^{2}-4t^{2}}{2\sqrt2 t^{2}}
=-\frac{1}{2\sqrt2}=-\frac{\sqrt2}{4}
\]
\end{solution}

\begin{problem}
数列 \(\{a_{n}\}\) 的通项公式 \(a_{n} = n\cos\frac{n\pi}{2} + 1\)，前 \(n\) 项和为 \(S_{n}\)，则 \(S_{2012} =\)\fillinblank{}.
\end{problem}

\begin{answer}
\(3018\)
\end{answer}

\begin{solution}
因为
\(\cos\frac{n\pi}{2}=0\)，\(-1\)，\(0\)，\(1\)
按 \(n\) 模 \(4\) 循环，故每四项中
\[
\begin{gathered}
(4k-3)\cos\frac{(4k-3)\pi}{2}+(4k-2)\cos\frac{(4k-2)\pi}{2}\\
+(4k-1)\cos\frac{(4k-1)\pi}{2}+4k\cos(2k\pi)=-(4k-2)+4k=2
\end{gathered}
\]
每四项的常数项 \(1\) 又贡献 \(4\)，所以每四项的和为 \(6\). 由于 \(2012=4\times503\)，
\[
S_{2012}=503\times6=3018
\]
\end{solution}

\begin{problem}
对于实数 \(a\) 和 \(b\)，定义运算“\(\ast\)”：\(a \ast b = \begin{cases}
a^{2} - ab, & a \leqslant b,\\
b^{2} - ab, & a > b.
\end{cases}\) 设 \(f(x) = (2x - 1) \ast (x - 1)\)，且关于 \(x\) 的方程 \(f(x) = m\) （\(m \in \R\)）恰有三个互不相等的实数根 \(x_{1}\)，\(x_{2}\)，\(x_{3}\)，则 \(x_{1}x_{2}x_{3}\) 的取值范围是\fillinblank{}.
\end{problem}

\begin{answer}
\(\left(\frac{1-\sqrt{3}}{16},0\right)\)
\end{answer}

\begin{solution}
比较运算中两个数的大小：
\[
2x-1\leqslant x-1\Longleftrightarrow x\leqslant0
\]
因此
\[
f(x)=\begin{cases}
2x^{2}-x,&x\leqslant0,\\
-x^{2}+x,&x>0
\end{cases}
\]
当 \(x\leqslant0\) 时，\(f(x)\) 从 \(+\infty\) 单调下降到 \(0\)，值域为 \([0,+\infty)\)；当 \(x>0\) 时，\(f(x)=\frac14-(x-\frac12)^{2}\)，最大值为 \(\frac14\).

当 \(m<0\) 时，左支无根，右支只有一个正根；当 \(m=0\) 时，根为 \(0,1\)，共有两个；当 \(0<m<\frac14\) 时，左支有一个负根，右支有两个正根；当 \(m=\frac14\) 时，右支只有一个根 \(x=\frac12\)，左支有一个负根；当 \(m>\frac14\) 时，只有左支的一个负根. 因此方程 \(f(x)=m\) 有三个互不相等的实根，当且仅当
\[
0<m<\frac14
\]
此时左支有一个负根，右支有两个正根. 设左支根为 \(x_{1}\)，右支两根为 \(x_{2}\)，\(x_{3}\)，则
\(2x_{1}^{2}-x_{1}-m=0\)，\(x^{2}-x+m=0\)
故
\(x_{1}=\frac{1-\sqrt{1+8m}}{4}\)，\(x_{2}x_{3}=m\)
从而
\[
x_{1}x_{2}x_{3}=\frac{m(1-\sqrt{1+8m})}{4}=:p(m)
\]
对 \(0<m<\frac14\)，
\[
p'(m)=\frac14\left(1-\frac{1+12m}{\sqrt{1+8m}}\right)<0
\]
其中不等号由 \((1+12m)^{2}-(1+8m)=16m+144m^{2}>0\) 得到. 因此 \(p(m)\) 严格递减，且
\(\lim_{m\to0^{+}}p(m)=0\)，\(\lim_{m\to(1/4)^{-}}p(m)=\frac{1-\sqrt3}{16}\)
端点均不取，故
\[
x_{1}x_{2}x_{3}\in\left(\frac{1-\sqrt3}{16},0\right)
\]
\end{solution}

\section{解答题}

\begin{problem}
受轿车在保修期内维修费等因素的影响，企业生产每辆轿车的利润与该轿车首次出现故障的时间有关. 某轿车制造厂生产甲，乙两种品牌轿车，保修期均为 \(2\) 年. 现从该厂已售出的两种品牌轿车中各随机抽取 \(50\) 辆，统计数据如下：

\begin{center}
\renewcommand{\arraystretch}{1.2}
\begin{tabular}{|c|c|c|c|c|c|}
\hline
品牌 & \multicolumn{3}{c|}{甲} & \multicolumn{2}{c|}{乙} \\
\hline
首次出现故障的时间 \(x\)（\(\text{年}\)） & \(0 < x \leqslant 1\) & \(1 < x \leqslant 2\) & \(x > 2\) & \(0 < x \leqslant 2\) & \(x > 2\) \\
\hline
轿车数量（\(\text{辆}\)） & 2 & 3 & 45 & 5 & 45 \\
\hline
每辆利润（\(\text{万元}\)） & 1 & 2 & 3 & 1.8 & 2.9 \\
\hline
\end{tabular}
\end{center}

将频率视为概率，解答下列问题：

\begin{enumerate}
\item 从该厂生产的甲品牌轿车中随机抽取一辆，求其首次出现故障发生在保修期内的概率；
\item 若该厂生产的轿车均能售出，记生产一辆甲品牌轿车的利润为 \(X_{1}\)，生产一辆乙品牌轿车的利润为 \(X_{2}\)，分别求 \(X_{1}\)，\(X_{2}\) 的分布列；
\item 该厂预计今后这两种品牌轿车销量相当，由于资金限制，只能生产其中一种品牌的新车. 若从经济效益的角度考虑，你认为应生产哪种品牌的轿车？请说明理由.
\end{enumerate}
\end{problem}

\begin{solution}
\begin{enumerate}
\item 甲品牌轿车首次出现故障发生在保修期内的数量为 \(2+3=5\) 辆，因此所求概率为
\[
\frac{5}{50}=\frac{1}{10}
\]

\item 由频率估计概率，\(X_{1}\) 的分布列为
\begin{center}
\begin{tabular}{c|ccc}
\hline
\(X_{1}\) & \(1\) & \(2\) & \(3\) \\
\hline
\(P\) & \(\dfrac{1}{25}\) & \(\dfrac{3}{50}\) & \(\dfrac{9}{10}\) \\
\hline
\end{tabular}
\end{center}

\(X_{2}\) 的分布列为
\begin{center}
\begin{tabular}{c|cc}
\hline
\(X_{2}\) & \(1.8\) & \(2.9\) \\
\hline
\(P\) & \(\dfrac{1}{10}\) & \(\dfrac{9}{10}\) \\
\hline
\end{tabular}
\end{center}

\item 两种品牌轿车的平均利润分别为
\[
E(X_{1})=1\times\frac{1}{25}+2\times\frac{3}{50}+3\times\frac{9}{10}
=\frac{143}{50}=2.86\text{（万元）}
\]
\[
E(X_{2})=1.8\times\frac{1}{10}+2.9\times\frac{9}{10}
=\frac{279}{100}=2.79\text{（万元）}
\]

因为 \(E(X_{1})>E(X_{2})\)，所以从经济效益角度应生产甲品牌轿车.
\end{enumerate}
\end{solution}

\begin{problem}
某同学在一次研究性学习中发现，以下五个式子的值都等于同一个常数：

\begin{circlelist}
\item \(\sin^{2}13^{\circ} + \cos^{2}17^{\circ} - \sin13^{\circ}\cos17^{\circ}\)；
\item \(\sin^{2}15^{\circ} + \cos^{2}15^{\circ} - \sin15^{\circ}\cos15^{\circ}\)；
\item \(\sin^{2}18^{\circ} + \cos^{2}12^{\circ} - \sin18^{\circ}\cos12^{\circ}\)；
\item \(\sin^{2}(-18^{\circ}) + \cos^{2}48^{\circ} - \sin(-18^{\circ})\cos48^{\circ}\)；
\item \(\sin^{2}(-25^{\circ}) + \cos^{2}55^{\circ} - \sin(-25^{\circ})\cos55^{\circ}\).
\end{circlelist}

\begin{enumerate}
\item 试从上述五个式子中选择一个，求出这个常数；
\item 根据（1）的计算结果，将该同学的发现推广为三角恒等式，并证明你的结论.
\end{enumerate}
\end{problem}

\begin{solution}
\begin{enumerate}
\item 选择第二个式子. 因为
\[
\sin^{2}15^{\circ}+\cos^{2}15^{\circ}-\sin15^{\circ}\cos15^{\circ}
=1-\frac{1}{2}\sin30^{\circ}
=1-\frac{1}{4}
=\frac{3}{4}
\]
所以这五个式子的公共值为 \(\dfrac{3}{4}\).

\item 当 \(\alpha+\beta=30^{\circ}\) 时，有三角恒等式
\[
\sin^{2}\alpha+\cos^{2}\beta-\sin\alpha\cos\beta=\frac{3}{4}
\]

事实上，由 \(\beta=30^{\circ}-\alpha\) 及和角公式，
\[
\cos\beta=\cos\left(30^{\circ}-\alpha\right)
=\frac{\sqrt{3}}{2}\cos\alpha+\frac{1}{2}\sin\alpha
\]
因此
\[
\sin^{2}\alpha+\cos^{2}\beta-\sin\alpha\cos\beta
=\sin^{2}\alpha+\left(\frac{\sqrt{3}}{2}\cos\alpha+\frac{1}{2}\sin\alpha\right)^{2}
-\sin\alpha\left(\frac{\sqrt{3}}{2}\cos\alpha+\frac{1}{2}\sin\alpha\right)
\]
展开并合并同类项，得
\[
\sin^{2}\alpha+\frac{3}{4}\cos^{2}\alpha+\frac{\sqrt{3}}{2}\sin\alpha\cos\alpha+\frac{1}{4}\sin^{2}\alpha
-\frac{\sqrt{3}}{2}\sin\alpha\cos\alpha-\frac{1}{2}\sin^{2}\alpha
=\frac{3}{4}\sin^{2}\alpha+\frac{3}{4}\cos^{2}\alpha
\]
所以
\[
\sin^{2}\alpha+\cos^{2}\beta-\sin\alpha\cos\beta
=\frac{3}{4}\left(\sin^{2}\alpha+\cos^{2}\alpha\right)
=\frac{3}{4}
\]
题中五组角度均满足 \(\alpha+\beta=30^{\circ}\)，故五个式子的值确实相同.
\end{enumerate}
\end{solution}

\begin{problem}
如图，在长方体 \(ABCD-A_{1}B_{1}C_{1}D_{1}\) 中，\(AA_{1} = AD = 1\)，\(E\) 为 \(CD\) 的中点.

\begin{enumerate}
\item 求证：\(B_{1}E \perp AD_{1}\)；
\item 在棱 \(AA_{1}\) 上是否存在一点 \(P\)，使得 \(DP \parallel\) 平面 \(B_{1}AE\)？若存在，求 \(AP\) 的长；若不存在，请说明理由；
\item 若二面角 \(A-B_{1}E-A_{1}\) 的大小为 \(30^{\circ}\)，求 \(AB\) 的长.
\end{enumerate}

\bitmapfigure[max width=0.56\linewidth]{img/2012/fujian_science/q18_fig1.png}
\end{problem}

\begin{solution}
\begin{enumerate}
\item 设 \(AB=b\)，其中 \(b>0\). 以 \(A\) 为原点，分别以 \(\overrightarrow{AB}\)，\(\overrightarrow{AD}\)，\(\overrightarrow{AA_{1}}\) 的方向为 \(x\) 轴、\(y\) 轴、\(z\) 轴的正方向建立空间直角坐标系，则 \(A(0,0,0)\)，\(B(b,0,0)\)，\(D(0,1,0)\)，\(A_{1}(0,0,1)\).
从而
\(B_{1}(b,0,1)\)，\(D_{1}(0,1,1)\)，\(E\left(\frac{b}{2},1,0\right)\).
于是 \(\overrightarrow{AD_{1}}=(0,1,1)\)，\(\overrightarrow{B_{1}E}=\left(-\frac{b}{2},1,-1\right)\).
\[
\overrightarrow{AD_{1}}\cdot\overrightarrow{B_{1}E}
=0\times\left(-\frac{b}{2}\right)+1\times1+1\times(-1)=0
\]
所以 \(B_{1}E\perp AD_{1}\).

\item 设 \(P(0,0,t)\)，其中 \(0\leqslant t\leqslant1\). 有 \(\overrightarrow{AB_{1}}=(b,0,1)\)，\(\overrightarrow{AE}=\left(\frac{b}{2},1,0\right)\).
取平面 \(B_{1}AE\) 的一个法向量 \(\bs{n}=\left(-1,\frac{b}{2},b\right)\)，因为 \(\bs{n}\cdot\overrightarrow{AB_{1}}=-b+b=0\)，\(\bs{n}\cdot\overrightarrow{AE}=-\frac{b}{2}+\frac{b}{2}=0\).
又 \(\overrightarrow{DP}=(0,-1,t)\).
所以
\[
\bs{n}\cdot\overrightarrow{DP}
=-\frac{b}{2}+bt=b\left(t-\frac{1}{2}\right)
\]
要使 \(DP\parallel\) 平面 \(B_{1}AE\)，应有 \(\bs{n}\cdot\overrightarrow{DP}=0\)，故 \(t=\dfrac{1}{2}\). 因此存在这样的点 \(P\)，且
\[
AP=t=\frac{1}{2}
\]

\item 平面 \(A_{1}B_{1}E\) 的一个法向量为 \(\bs{n}_{2}=(0,1,1)\)，平面 \(B_{1}AE\) 的法向量为 \(\bs{n}_{1}=\left(-1,\frac{b}{2},b\right)\). 因此两平面夹角 \(\theta\) 满足
\[
\cos\theta
=\frac{|\bs{n}_{1}\cdot\bs{n}_{2}|}{|\bs{n}_{1}|\,|\bs{n}_{2}|}
=\frac{\frac{3b}{2}}{\sqrt{1+\frac{5b^{2}}{4}}\sqrt{2}}
\]
由 \(\theta=30^{\circ}\)，得
\[
\frac{9b^{2}}{8+10b^{2}}=\frac{3}{4}
\]
解得 \(b^{2}=4\). 因为 \(b>0\)，所以 \(AB=b=2\).
\end{enumerate}
\end{solution}

\begin{problem}
如图，椭圆 \(E: \frac{x^{2}}{a^{2}} + \frac{y^{2}}{b^{2}} = 1\)（\(a > b > 0\)）的左焦点为 \(F_{1}\)，右焦点为 \(F_{2}\)，离心率 \(\varepsilon = \frac{1}{2}\)，过 \(F_{1}\) 的直线交椭圆于 \(A\)，\(B\) 两点，且 \(\triangle ABF_{2}\) 的周长为 \(8\).

\begin{enumerate}
\item 求椭圆 \(E\) 的方程；
\item 设动直线 \(l: y = kx + m\) 与椭圆 \(E\) 有且只有一个公共点 \(P\)，且与直线 \(x = 4\) 相交于点 \(Q\). 试探究：在坐标平面内是否存在定点 \(M\)，使得以 \(PQ\) 为直径的圆恒过点 \(M\)？若存在，求出点 \(M\) 的坐标；若不存在，请说明理由.
\end{enumerate}

\FigureLayoutDeclare{img/2012/fujian_science/q19_fig1.png}{9.5cm}{89bp 86bp 71bp 38bp}
\begin{center}
\bitmapfigure[float=inline,max width=0.48\linewidth]{img/2012/fujian_science/q19_fig1.png}
\end{center}
\end{problem}

\begin{solution}
\begin{enumerate}
\item 设椭圆的半焦距为 \(c\). 由于 \(F_{1}\) 在椭圆内部且直线 \(AB\) 经过 \(F_{1}\)，所以 \(F_{1}\) 在线段 \(AB\) 上，从而
\[
8=AB+AF_{2}+BF_{2}
=(AF_{1}+BF_{1})+AF_{2}+BF_{2}
=(AF_{1}+AF_{2})+(BF_{1}+BF_{2})
=4a
\]
故 \(a=2\). 又由离心率
\[
\varepsilon=\frac{c}{a}=\frac{1}{2}
\]
得 \(c=1\)，于是
\[
b^{2}=a^{2}-c^{2}=4-1=3
\]
因此椭圆 \(E\) 的方程为
\[
\frac{x^{2}}{4}+\frac{y^{2}}{3}=1
\]

\item 设切点 \(P=(x_{0},y_{0})\)，并设 \(Q=(4,q)\). 若 \(y_{0}=0\)，椭圆在 \(P\) 处的切线为 \(x=2\) 或 \(x=-2\)，均与直线 \(x=4\) 平行，故题设中的切线必有 \(y_{0}\neq0\).

椭圆在 \(P\) 处的切线方程为
\[
\frac{x_{0}x}{4}+\frac{y_{0}y}{3}=1
\]
将 \(Q=(4,q)\) 代入，得
\[
x_{0}+\frac{y_{0}q}{3}=1
\]
即
\[
y_{0}q=3(1-x_{0})
\]
取定点 \(M=(1,0)\)，则
\(\overrightarrow{PM}=(1-x_{0},-y_{0})\)，\(\overrightarrow{QM}=(-3,-q)\)
于是
\[
\overrightarrow{PM}\cdot\overrightarrow{QM}
=-3(1-x_{0})+y_{0}q=0
\]
所以 \(PM\perp QM\)，从而 \(\angle PMQ=90^{\circ}\). 由圆的直径所对的圆周角为直角，点 \(M\) 在以 \(PQ\) 为直径的圆上. 因此存在定点 \(M(1,0)\).
\end{enumerate}
\end{solution}

\FloatBarrier

\begin{problem}
已知函数 \(f(x) = \e^{x} + ax^{2} - \e x\)，\(a \in \R\).

\begin{enumerate}
\item 若曲线 \(y = f(x)\) 在点 \((1,f(1))\) 处的切线平行于 \(x\) 轴，求函数 \(f(x)\) 的单调区间；
\item 试确定 \(a\) 的取值范围，使得曲线 \(y = f(x)\) 上存在唯一的点 \(P\)，曲线在该点处的切线与曲线只有一个公共点 \(P\).
\end{enumerate}
\end{problem}

\begin{solution}
\begin{enumerate}
\item \(f'(x)=\e^{x}+2ax-\e\). 由题意，曲线在 \(x=1\) 处的切线平行于 \(x\) 轴，所以
\[
f'(1)=\e+2a-\e=2a=0
\]
解得 \(a=0\). 此时
\[
f'(x)=\e^{x}-\e
\]
当 \(x<1\) 时，\(\e^{x}<\e\)，故 \(f'(x)<0\)；当 \(x>1\) 时，\(\e^{x}>\e\)，故 \(f'(x)>0\). 因此 \(f(x)\) 的单调递减区间为 \((-\infty,1)\)，单调递增区间为 \((1,+\infty)\).

\item 设 \(P=(t,f(t))\). 曲线在 \(P\) 处的切线方程为
\[
y=f'(t)(x-t)+f(t)
\]
记曲线与该切线的纵坐标之差为
\[
G_{t}(x)=f(x)-f'(t)(x-t)-f(t)
\]
令 \(x=t+h\)，则
\[
G_{t}(t+h)=\e^{t}\left(\e^{h}-1-h\right)+ah^{2}
=h^{2}\left(a+\e^{t}\frac{\e^{h}-1-h}{h^{2}}\right)
\]
当 \(h=0\) 时，\(G_{t}(t)=0\). 对 \(h\neq0\)，记
\[
r(h)=\frac{\e^{h}-1-h}{h^{2}}
\]
由 \(\e^{h}>1+h\)（\(h\neq0\)），知 \(r(h)>0\). 再令
\[
\phi(h)=\e^{h}-1-h-\frac{h^{2}}{2}
\]
由于 \(\phi(0)=\phi'(0)=0\)，且
\[
\phi''(h)=\e^{h}-1
\]
当 \(h>0\) 时，\(\phi''(h)>0\)，故 \(\phi'\) 在 \((0,+\infty)\) 上递增，从而 \(\phi'(h)>\phi'(0)=0\)，进而 \(\phi(h)>\phi(0)=0\)；当 \(h<0\) 时，\(\phi''(h)<0\)，故 \(\phi'\) 在 \((-\infty,0)\) 上递减，所以 \(\phi'(h)>\phi'(0)=0\)，进而 \(\phi(h)<\phi(0)=0\). 因此
\(\phi(h)<0\)，\((h<0)\)，\(\phi(h)>0\)，\((h>0)\)
于是
\(0<r(h)<\frac{1}{2}\)，\((h<0)\)，\(r(h)>\frac{1}{2}\)，\((h>0)\)
此外，\(r(h)\) 在 \(h\neq0\) 时连续，且
\(\lim_{h\to-\infty}r(h)=0\)，\(\lim_{h\to0}r(h)=\frac{1}{2}\)，\(\lim_{h\to+\infty}r(h)=+\infty\)

若 \(a\geqslant0\)，则对任意 \(t\) 及 \(h\neq0\)，
\[
G_{t}(t+h)=h^{2}\left(a+\e^{t}r(h)\right)>0
\]
此时每个点处的切线都只与曲线交于该点，不满足“存在唯一的点 \(P\)”.

若 \(a<0\)，令
\[
t_{0}=\ln(-2a)
\]
则 \(\e^{t_{0}}=-2a\). 当 \(t=t_{0}\) 时，
\[
G_{t_{0}}(t_{0}+h)
=\e^{t_{0}}h^{2}\left(r(h)-\frac{1}{2}\right)
\]
由 \(r(h)<\dfrac{1}{2}\)（\(h<0\)）及 \(r(h)>\dfrac{1}{2}\)（\(h>0\)），知 \(G_{t_{0}}(t_{0}+h)\neq0\)（\(h\neq0\)）. 因此在 \(P=(t_{0},f(t_{0}))\) 处，切线与曲线只有一个公共点.

当 \(t\neq t_{0}\) 时，\(c=-a\e^{-t}>0\) 且 \(c\neq\dfrac{1}{2}\). 由 \(r(h)\) 在负半轴上的取值覆盖 \((0,\dfrac{1}{2})\)，在正半轴上的取值覆盖 \((\dfrac{1}{2},+\infty)\)，可取 \(h\neq0\) 使 \(r(h)=c\). 于是 \(G_{t}(t+h)=0\)，该切线还与曲线有另一个公共点. 所以只有 \(t=t_{0}\) 对应的点具有题设性质，参数范围为
\[
a<0
\]
\end{enumerate}
\end{solution}

\section{选考题：解答题}

\begin{problem}
设曲线 \(2x^{2} + 2xy + y^{2} = 1\) 在矩阵 \(\mathbf{A} = \begin{pmatrix} a & 0 \\ b & 1 \end{pmatrix}\) （\(a > 0\)）对应的变换作用下得到的曲线为 \(x^{2} + y^{2} = 1\).

\begin{enumerate}
\item 求实数 \(a\)，\(b\) 的值；
\item 求 \(A^{2}\) 的逆矩阵.
\end{enumerate}
\end{problem}

\begin{solution}
\begin{enumerate}
\item 设原曲线上任意一点 \(P(x,y)\) 在矩阵 \(\mathbf{A}\) 对应的变换下的像为 \(P'(x',y')\)，则
\[
\begin{pmatrix}x'\\y'\end{pmatrix}
=\begin{pmatrix}a&0\\b&1\end{pmatrix}
\begin{pmatrix}x\\y\end{pmatrix}
\]
所以
\(x'=ax\)，\(y'=bx+y\)
像曲线为 \(x'^{2}+y'^{2}=1\)，故原曲线上的点满足
\[
a^{2}x^{2}+(bx+y)^{2}=1
\]
展开得
\[
(a^{2}+b^{2})x^{2}+2bxy+y^{2}=1
\]
与原曲线方程 \(2x^{2}+2xy+y^{2}=1\) 比较系数，得到
\(a^{2}+b^{2}=2\)，\(2b=2\)
因此 \(b=1\)，\(a^{2}=1\). 由 \(a>0\)，得
\(a=1\)，\(b=1\)

\item 由（1）得
\[
\mathbf{A}=\begin{pmatrix}1&0\\1&1\end{pmatrix},\qquad
\mathbf{A}^{2}
=\begin{pmatrix}1&0\\1&1\end{pmatrix}
\begin{pmatrix}1&0\\1&1\end{pmatrix}
=\begin{pmatrix}1&0\\2&1\end{pmatrix}
\]
因为
\[
\begin{pmatrix}1&0\\2&1\end{pmatrix}
\begin{pmatrix}1&0\\-2&1\end{pmatrix}
=\begin{pmatrix}1&0\\0&1\end{pmatrix}
\]
所以
\[
(\mathbf{A}^{2})^{-1}=\begin{pmatrix}1&0\\-2&1\end{pmatrix}
\]
\end{enumerate}
\end{solution}

\begin{problem}
在平面直角坐标系中，以坐标原点 \(O\) 为极点，\(x\) 轴的正半轴为极轴建立极坐标系. 已知直线 \(l\) 上两点 \(M\)，\(N\) 的极坐标分别为 \((2,0)\)，\(\left(\frac{2\sqrt{3}}{3},\frac{\pi}{2}\right)\)，圆 \(C\) 的参数方程为 \(\begin{cases}
x = 2 + 2\cos\theta,\\
y = -\sqrt{3} + 2\sin\theta,
\end{cases}\)（\(\theta\) 为参数）.

\begin{enumerate}
\item 设 \(P\) 为线段 \(MN\) 的中点，求直线 \(OP\) 的平面直角坐标方程；
\item 判断直线 \(l\) 与圆 \(C\) 的位置关系.
\end{enumerate}
\end{problem}

\begin{solution}
\begin{enumerate}
\item 极坐标与平面直角坐标的转换关系为 \(x=\rho\cos\theta\)，\(y=\rho\sin\theta\). 因此
\(M=(2,0)\)，\(N=\left(0,\frac{2\sqrt{3}}{3}\right)\)
线段 \(MN\) 的中点为
\[
P=\left(1,\frac{\sqrt{3}}{3}\right)
\]
直线 \(OP\) 经过原点，斜率为 \(\dfrac{\sqrt{3}}{3}\)，所以其方程为
\[
y=\frac{\sqrt{3}}{3}x
\]
即
\[
x-\sqrt{3}y=0
\]

\item 直线 \(l\) 经过 \(M\) 和 \(N\)，故
\[
l:\ x+\sqrt{3}y=2
\]
由圆 \(C\) 的参数方程可知，圆心为 \(C_{0}=(2,-\sqrt{3})\)，半径为 \(2\). 圆心到直线 \(l\) 的距离为
\[
d=\frac{|2+\sqrt{3}(-\sqrt{3})-2|}{\sqrt{1+3}}
=\frac{3}{2}<2
\]
所以直线 \(l\) 与圆 \(C\) 相交.
\end{enumerate}
\end{solution}

\begin{problem}
已知函数 \(f(x) = m - |x - 2|\)，\(m \in \R\). 且 \(f(x + 2) \geqslant 0\) 的解集为 \([-1,1]\).

\begin{enumerate}
\item 求 \(m\) 的值；
\item 若 \(a\)，\(b\)，\(c \in \R^{+}\)，且 \(\frac{1}{a} + \frac{1}{2b} + \frac{1}{3c} = m\)，求证：\(a + 2b + 3c \geqslant 9\).
\end{enumerate}
\end{problem}

\begin{solution}
\begin{enumerate}
\item
\[
f(x+2)=m-|x|
\]
所以不等式 \(f(x+2)\geqslant0\) 等价于
\[
|x|\leqslant m
\]
要使解集为非空的区间 \([-1,1]\)，必须有 \(m\geqslant0\)，且此时 \(|x|\leqslant m\) 的解集为 \([-m,m]\). 因此
\[
m=1
\]

\item 由（1）知
\[
\frac{1}{a}+\frac{1}{2b}+\frac{1}{3c}=1
\]
因为 \(a\)，\(b\)，\(c>0\)，由柯西不等式
\[
\left(a+2b+3c\right)\left(\frac{1}{a}+\frac{1}{2b}+\frac{1}{3c}\right)
\geqslant
\left(\sqrt{a}\cdot\frac{1}{\sqrt{a}}
+\sqrt{2b}\cdot\frac{1}{\sqrt{2b}}
+\sqrt{3c}\cdot\frac{1}{\sqrt{3c}}\right)^{2}
=9
\]
代入 \(\dfrac{1}{a}+\dfrac{1}{2b}+\dfrac{1}{3c}=1\)，得
\[
a+2b+3c\geqslant9
\]
\end{enumerate}
\end{solution}
